% Generated draft: requires mathematical and source review.
\documentclass[11pt]{article}
\usepackage[margin=25mm]{geometry}
\usepackage{fontspec}
\setmainfont{Latin Modern Roman}
\usepackage{amsmath,amssymb,mathtools,unicode-math}
\setmathfont{Latin Modern Math}
\usepackage{enumitem,needspace,xurl,hyperref}
\setlength{\parindent}{0pt}
\setlength{\parskip}{5pt}
\setlength{\emergencystretch}{3em}
\newcommand{\N}{\mathbb{N}}
\newcommand{\Z}{\mathbb{Z}}
\newcommand{\Q}{\mathbb{Q}}
\newcommand{\R}{\mathbb{R}}
\newcommand{\C}{\mathbb{C}}
\newcommand{\F}{\mathbb{F}}
\newcommand{\A}{\mathbb{A}}
\newcommand{\PP}{\mathbb{P}}
\newcommand{\E}{\mathbb{E}}
\newcommand{\eps}{\varepsilon}
\newcommand{\dd}{\,\mathrm{d}}
\begin{document}
\section*{UnsolvedMath ID: 7000015}
\textbf{Source problem code:} catalogue entry
\subsection{Definitions and mathematical statement}
Let $P\subset\R^3$ be a full-dimensional compact convex polyhedron and let $S=\partial P$ with its intrinsic path metric $d_S$, measuring shortest surface-path length. A genuine edge is a one-dimensional face between distinct maximal planar facets; a diagonal added to triangulate one facet is not a genuine edge.

The source asks for a reasonably simple algorithm that recovers these edges using intrinsic data for $(S,d_S)$. It does not prescribe the input representation, but possible finite data would be Euclidean triangle gluings describing the intrinsic polyhedral metric without supplying its spatial realization. The requested output must identify the arcs on $S$ corresponding to genuine edges, not merely any triangulation of the metric.

Intrinsic cone curvature detects genuine vertices through their angle deficits. By contrast a nonvertex point in the interior of an edge has an intrinsically flat neighborhood: the adjacent faces unfold to a planar patch. Thus a nontrivial spatial dihedral angle is not itself a local intrinsic metric singularity. A global reconstruction argument and a direct local edge test are different methods.

The supplied summary invokes Aleksandrov-type reconstruction/uniqueness, which explains global determinacy under its appropriate convex-metric hypotheses, but does not specify a simple recovery algorithm. Whether an algorithm may reconstruct the spatial polyhedron as an intermediate step, what ``reasonably simple'' means, which exact coordinate model is used, and how coplanarity/degeneracies are decided are all absent. The prompt is therefore a request for a useful intrinsic procedure rather than a precisely bounded complexity conjecture.
\subsection{Short English statement}
Recover the genuine facet edges of a convex polyhedron from its intrinsic surface metric by a simple procedure. Vertices are intrinsic curvature points, while an ordinary edge point is locally flat, so the required edge information is global.
\subsection{Sources}
\begin{itemize}
\item Local UnsolvedMath registry, permanent ID 7000015, exact source code catalogue entry, ``Geometry of Curves and Surfaces---Problem 3.2''; proposed by Mohammad Ghomi, 2019. Registry status open, classification OPEN-TRIAGE.
\item Mohammad Ghomi, \emph{Open Problems in Geometry of Curves and Surfaces} (2019), Problem 3.2, PDF page 11, \url{https://people.math.gatech.edu/~ghomi/Papers/op.pdf}. Supplied normalized pdf-text extraction accessed July 29, 2026. No PDF was accessed here.
\item The supplied summary mentions Aleksandrov reconstruction and source unfolding without precise references and says no canonical simple intrinsic edge-detection algorithm was verified. Its description of a nontrivial dihedral angle is extrinsic and is not used as a local intrinsic detection criterion.
\item Provenance flags current status and statement-reuse terms NEEDS\_REVIEW; the source is publicly author-hosted.
\end{itemize}
\end{document}
